\documentclass[11pt]{article}
\usepackage{amsmath}
\usepackage{amssymb}
\usepackage{geometry}
\geometry{margin=1in}

\title{Plan 02-01: Atmospheric Drag Acceleration}
\author{J2+J3+Drag Secular Perturbation Theory}
\date{2026-04-02 — Phase 2}

\begin{document}
\maketitle

\section{Objective}
Express atmospheric drag force and acceleration in terms that can be integrated with orbital element framework.

\section{Drag Force}

Atmospheric drag on a satellite:
\begin{equation}
\vec{F}_D = -\frac{1}{2}\rho v^2 C_D A \,\hat{v}
\end{equation}

where:
\begin{itemize}
    \item $\rho$ — atmospheric density [kg/m³]
    \item $v$ — orbital velocity magnitude [m/s]
    \item $C_D$ — drag coefficient [dimensionless], typically 2.0-2.5
    \item $A$ — cross-sectional area [m²]
    \item $\hat{v}$ — velocity unit vector (opposite to motion)
\end{itemize}

\section{Drag Acceleration}

Dividing by satellite mass $m$:
\begin{equation}
\boxed{\vec{a}_D = -\frac{1}{2}\rho v^2 B\,\hat{v}}
\label{eq:drag_accel}
\end{equation}

where the \textbf{ballistic coefficient} is:
\begin{equation}
B = \frac{C_D A}{m} \quad [\text{m}^2/\text{kg}]
\end{equation}

Typical values: $B \approx 0.01$-$0.1$ m²/kg depending on satellite design.

\section{Exponential Atmosphere Model}

Atmospheric density decreases exponentially with altitude:
\begin{equation}
\boxed{\rho(h) = \rho_0 \exp\left(-\frac{h}{H}\right)}
\label{eq:exponential_density}
\end{equation}

where:
\begin{itemize}
    \item $h = r - R_E$ — altitude above Earth's surface [km]
    \item $H$ — atmospheric scale height [km], typically 60-80 km (altitude-dependent)
    \item $\rho_0$ — reference density at altitude $h=0$ or specific reference altitude
\end{itemize}

\section{Orbital Velocity}

From the vis-viva equation:
\begin{equation}
v^2 = \mu\left(\frac{2}{r} - \frac{1}{a}\right)
\label{eq:vis_viva}
\end{equation}

For circular orbits ($r = a$):
\begin{equation}
v_{\text{circular}} = \sqrt{\frac{\mu}{a}}
\end{equation}

For near-circular orbits with small eccentricity:
\begin{equation}
v \approx \sqrt{\frac{\mu}{a}}\left[1 + O(e)\right]
\end{equation}

\section{Drag Acceleration Magnitude}

Substituting Eqs.~\eqref{eq:exponential_density} and \eqref{eq:vis_viva} into \eqref{eq:drag_accel}:

\begin{equation}
a_D = \frac{1}{2}\rho_0 \exp\left(-\frac{r - R_E}{H}\right) \cdot \mu\left(\frac{2}{r} - \frac{1}{a}\right) \cdot B
\end{equation}

For circular orbits ($r = a$):
\begin{equation}
\boxed{a_{D,\text{circular}} = \frac{1}{2}\rho_0 B \frac{\mu}{a} \exp\left(-\frac{a - R_E}{H}\right)}
\label{eq:drag_circular}
\end{equation}

\section{Dimensional Analysis}

\begin{align}
[a_D] &= [\rho][v^2][B]\\
&= (\text{kg/m}^3)(\text{m}^2/\text{s}^2)(\text{m}^2/\text{kg})\\
&= \text{m/s}^2 \quad \checkmark
\end{align}

Correct dimensions for acceleration.

\section{Direction: Tangential Component}

Drag acts opposite to velocity. In the orbital plane:
\begin{itemize}
    \item \textbf{Radial component:} $a_R \approx 0$ (velocity is primarily tangential)
    \item \textbf{Tangential component:} $a_T = -a_D$ (opposes motion)
    \item \textbf{Normal component:} $a_N = 0$ (drag is in orbital plane)
\end{itemize}

This will be formalized in Plan 02-03 using Gauss equations.

\section{Success Criteria}
\begin{itemize}
    \item[$\checkmark$] Drag acceleration vector expressed
    \item[$\checkmark$] Ballistic coefficient $B$ defined
    \item[$\checkmark$] Exponential density model stated
    \item[$\checkmark$] Dimensional analysis passed
\end{itemize}

\end{document}
